\chapter{EDA Lab I: The EDA Playbook in Code}
\companion{Krish Naik: Exploratory Data Analysis with pandas (search)}{\yts{exploratory+data+analysis+pandas+python+complete+tutorial}}

EDA carries heavy weight in the InfoEdge interviews: Round 3 is literally ``clean the data, analyse it, model it, test it'', and the
successful candidate's first moves (spotting salary leakage, checking ranges before calling anything an outlier) were EDA judgements. This part of the
book is hands-on. It comes with a \textbf{Colab-ready notebook} that runs every analysis shown here on real data:
\begin{tcolorbox}[colback=accentlight,colframe=accent,boxrule=0.5pt,arc=2pt]
\small\faLaptopCode\ \textbf{Notebook:} \code{InfoEdge\_Interview\_Companion/eda\_lab/EDA\_Lab.ipynb} in your repository (branch
\code{claude/admiring-edison-vd2vbx}). Open it in Google Colab (File $\to$ Open notebook $\to$ GitHub, or upload the file) and run all cells. The real
datasets (HR job change, Titanic, Pima diabetes) are read directly from your public GitHub repository. For the MBA placement case, upload Kaggle's
\code{Placement\_Data\_Full\_Class.csv} to the Colab session; without it the notebook builds a simulated replica with the same columns.
\end{tcolorbox}
Chapter 41 gives the method and a cookbook of pandas/NumPy recipes, each run on the real HR job-change dataset (19{,}158 data-science candidates; target
$=1$ if they are looking for a job change: very close to Naukri data). Chapters 42--44 are complete case studies.

\section{The 12-step EDA workflow}
Say the step, run the code, read the output, state the decision. Every step ends in ``\emph{because I see X, I will do Y}''.
\begin{center}\small
\begin{tabularx}{\linewidth}{r l X}
\toprule
& Step & Question it answers\\
\midrule
1 & Frame & What is one row? What is the target? When is a prediction made?\\
2 & Load and look & Shape, \code{head}, types, memory; does it look like the description?\\
3 & Types and parsing & Numbers stored as text? Codes like ``>20'', ``never''? Dates?\\
4 & Data quality & Duplicates, impossible values, malformed labels, constant columns, disguised missing values (0, $-1$, 999)\\
5 & Missing values & How much, where, together with what, and does missingness relate to the target?\\
6 & Target & Balance; base rate; metric implications\\
7 & Univariate numeric & Centre, spread, skew, outliers (genuine or error?)\\
8 & Univariate categorical & Levels, rare levels, cardinality\\
9 & Bivariate with target & Which features separate the classes, linearly or not?\\
10 & Feature--feature & Correlations, redundancy, multicollinearity\\
11 & Multivariate & Interactions and segments (does an effect depend on another variable?)\\
12 & Leakage and summary & Anything ``too good''? Write the findings and the modelling plan\\
\bottomrule
\end{tabularx}
\end{center}

\section{The reusable toolkit}
Define these once at the top of any interview notebook (they are cell 3 of \code{EDA\_Lab.ipynb}).
\begin{lstlisting}
def overview(df):
    # One row per column: type, missing %, number of unique values, an example value.
    return pd.DataFrame({
        "dtype": df.dtypes.astype(str),
        "missing": df.isna().sum(),
        "missing_%": (df.isna().mean() * 100).round(1),
        "n_unique": df.nunique(),
        "example": df.apply(lambda s: s.dropna().iloc[0] if s.notna().any() else np.nan)})

def num_summary(df):
    # describe() + skew, kurtosis, IQR-outlier count and zero count for numeric columns.
    num = df.select_dtypes("number")
    q1, q3 = num.quantile(0.25), num.quantile(0.75); iqr = q3 - q1
    res = num.describe().T
    res["skew"], res["kurt"] = num.skew(), num.kurt()
    res["n_iqr_outliers"] = ((num < q1 - 1.5 * iqr) | (num > q3 + 1.5 * iqr)).sum()
    res["n_zeros"] = (num == 0).sum()
    return res.round(2)

def rate_by(df, col, target, bins=None, q=None):
    # Target rate and count per category, or per bin of a numeric column.
    key = df[col]
    if bins is not None: key = pd.cut(key, bins)
    elif q is not None:  key = pd.qcut(key, q, duplicates="drop")
    key = key.astype(object).where(key.notna(), "MISSING")
    return df.groupby(key, observed=True)[target].agg(rate="mean", n="size").round(3)

def missing_vs_target(df, target):
    # Is a column's missingness related to the target? (informative missingness / leakage)
    rows = []
    for c in df.columns.drop(target):
        m = df[c].isna()
        if 0 < m.sum() < len(df):
            rows.append((c, m.mean() * 100, df.loc[m, target].mean(), df.loc[~m, target].mean()))
    return pd.DataFrame(rows, columns=["column", "missing_%", "rate_if_missing", "rate_if_present"]).round(3)

def point_biserial(df, num_cols, target):          # numeric feature vs 0/1 target
    rows = [(c, *stats.pointbiserialr(df[[c, target]].dropna()[target],
                                      df[[c, target]].dropna()[c])) for c in num_cols]
    return pd.DataFrame(rows, columns=["feature", "r_pb", "p_value"]).round(4)

def cramers_v(x, y):                               # categorical vs categorical, 0..1
    t = pd.crosstab(x, y); chi2 = stats.chi2_contingency(t, correction=False)[0]
    return np.sqrt(chi2 / (t.values.sum() * (min(t.shape) - 1)))

def vif(df_num):                                   # VIF = 1 / (1 - R^2_j), via numpy least squares
    X = df_num.dropna().values; out = {}
    for j, c in enumerate(df_num.columns):
        y = X[:, j]; A = np.column_stack([np.ones(len(X)), np.delete(X, j, axis=1)])
        coef, *_ = np.linalg.lstsq(A, y, rcond=None)
        r2 = 1 - ((y - A @ coef) ** 2).sum() / ((y - y.mean()) ** 2).sum()
        out[c] = 1 / (1 - r2)
    return pd.Series(out, name="VIF").round(2)
\end{lstlisting}

\section{The cookbook: 25 recipes, run on real data}
All outputs below come from running the code on \code{aug\_train.csv} (HR job change). Load: \code{df = pd.read\_csv("Data\_gathering/dataset/aug\_train.csv")}.

\subsection{Steps 2--4: structure, types, quality}
\textbf{R1. First look.}
\begin{lstlisting}
print(df.shape); print(df.dtypes.value_counts()); print(df.memory_usage(deep=True).sum() / 1e6, "MB")
\end{lstlisting}
\outlabel
\begin{lstlisting}[style=out]
(19158, 14)    10 text columns, 2 int64, 2 float64    12.39 MB
\end{lstlisting}
\textbf{R5. Duplicates} (exact, and ignoring the ID column, which hides duplicates):
\begin{lstlisting}
df.duplicated().sum(), df.drop(columns="enrollee_id").duplicated().sum(), df["enrollee_id"].is_unique
\end{lstlisting}
\outlabel
\begin{lstlisting}[style=out]
(0, 49, True)      # 49 people identical in every field except their ID
\end{lstlisting}
Decision: 49 of 19{,}158 is negligible; check whether they are the same person registered twice (keep one) or coincidence (keep all). In a group split they must
stay together.

\textbf{R6. Cardinality and constant columns:} \code{df.nunique().sort\_values()} $\to$ \code{relevent\_experience 2, target 2, enrolled\_university 3, gender 3, \dots, city 123,
enrollee\_id 19158}. ID-like columns (unique per row) are dropped; constant columns carry no information.

\textbf{R8. Parsing numbers stored as text.}
\begin{lstlisting}
exp = df["experience"].replace({">20": "21", "<1": "0"})
df["exp_num"] = pd.to_numeric(exp, errors="coerce")       # unparsable -> NaN, never crash
df["exp_num"].isna().sum()                                # 65 = exactly the originally missing values
\end{lstlisting}
\textbf{R7. Value counts with percentages, including missing.}
\begin{lstlisting}
vc = df["education_level"].value_counts(dropna=False)
pd.concat([vc, (vc / len(df) * 100).round(1)], axis=1, keys=["n", "%"])
\end{lstlisting}
\outlabel
\begin{lstlisting}[style=out]
                     n     %
Graduate         11598  60.5
Masters           4361  22.8
High School       2017  10.5
NaN                460   2.4
Phd                414   2.2
Primary School     308   1.6
\end{lstlisting}
\textbf{R9. Grouping rare levels} (cities with $<$0.5\% of rows):
\begin{lstlisting}
freq = df["city"].value_counts(normalize=True)
rare = freq[freq < 0.005].index
df["city_grp"] = df["city"].where(~df["city"].isin(rare), "Other")   # 89 rare cities -> "Other"; 35 levels left
\end{lstlisting}

\subsection{Step 5: missing values}
\textbf{R2. Missing table, sorted.}
\begin{lstlisting}
df.isna().agg(["sum", "mean"]).T.sort_values("sum", ascending=False).head()
\end{lstlisting}
\outlabel
\begin{lstlisting}[style=out]
                     sum   mean
company_type      6140  0.320
company_size      5938  0.310
gender            4508  0.235
major_discipline  2813  0.147
education_level    460  0.024
\end{lstlisting}
\textbf{R3. How many fields are missing per row?} \code{df.isna().sum(axis=1).value\_counts().sort\_index()} $\to$ \code{\{0: 8955, 1: 3718, 2: 3654, 3: 1953, 4: 628,
5: 176, 6: 62, 7: 12\}}. Only 47\% of rows are complete: dropping rows with any NaN would discard more than half the data.

\textbf{R4. Do fields go missing together?}
\begin{lstlisting}
m = df[["company_size", "company_type", "gender"]].isna()
m.corr()                                    # correlation of the missingness indicators
pd.crosstab(m.company_size, m.company_type)
\end{lstlisting}
\outlabel
\begin{lstlisting}[style=out]
              company_size  company_type  gender
company_size          1.00          0.84    0.06
company_type          0.84          1.00    0.07
company_type  False  True
company_size
False         12440    780
True            578   5360
\end{lstlisting}
Reading: company size and type are usually missing \emph{together} (5{,}360 rows), suggesting one cause: the candidate has no current company.

\textbf{R23. Missingness vs target and a missing indicator.}
\begin{lstlisting}
df["company_size_missing"] = df["company_size"].isna().astype(int)
df.groupby("company_size_missing")["target"].mean()           # {0: 0.179, 1: 0.406}
\end{lstlisting}
Blank company size more than doubles the job-change rate: informative missingness. Keep the indicator; never fill with the mode.

\subsection{Steps 6--8: target and univariate}
\textbf{R10. Percentiles that reveal tails.}
\begin{lstlisting}
df["training_hours"].describe(percentiles=[.01, .05, .5, .95, .99])
\end{lstlisting}
\outlabel
\begin{lstlisting}[style=out]
mean 65.4  std 60.1  min 1  1% 3  5% 7  50% 47  95% 188  99% 302  max 336
\end{lstlisting}
Mean (65) far above median (47): right skew.
\textbf{R11. Skew of every numeric column:} \code{df.select\_dtypes("number").skew()} $\to$ \code{city\_development\_index $-1.00$, training\_hours 1.82, exp\_num 0.41}.
\textbf{R25. Does a log fix it?} \code{np.log1p(df.training\_hours).skew()} $\to$ $-0.35$ (from 1.82).

\textbf{R12. IQR outlier flags.}
\begin{lstlisting}
q1, q3 = df.training_hours.quantile([.25, .75]); iqr = q3 - q1
flag = (df.training_hours < q1 - 1.5 * iqr) | (df.training_hours > q3 + 1.5 * iqr)
q3 + 1.5 * iqr, flag.sum(), flag.mean() * 100            # (185.5, 984, 5.1 %)
\end{lstlisting}
\textbf{R13. z-score flags with NumPy:} \code{z = (x - x.mean()) / x.std(); (np.abs(z) > 3).sum()} $\to$ 450.
Reading: on a right-skewed variable both rules flag hundreds of \emph{genuine} heavy learners. They are not errors: keep them; log-transform or cap for linear models.

\subsection{Step 9: features vs target}
\textbf{R14. Rate with counts per category.}
\begin{lstlisting}
df.groupby("relevent_experience")["target"].agg(rate="mean", n="size")
\end{lstlisting}
\outlabel
\begin{lstlisting}[style=out]
                          rate      n
Has relevent experience  0.215  13792
No relevent experience   0.338   5366
\end{lstlisting}
\textbf{R15. Crosstab normalised by row:}
\begin{lstlisting}
pd.crosstab(df["enrolled_university"], df["target"], normalize="index")
\end{lstlisting}
$\to$ job-change share: full-time course 0.381, part-time 0.252, no enrolment 0.211.

\textbf{R16. Numeric feature in quantile bins vs target.}
\begin{lstlisting}
df.groupby(pd.qcut(df["training_hours"], 4))["target"].mean()
\end{lstlisting}
\outlabel
\begin{lstlisting}[style=out]
(0.999, 23.0]    0.250
(23.0, 47.0]     0.258
(47.0, 88.0]     0.253
(88.0, 336.0]    0.237
\end{lstlisting}
Flat: training hours carry almost no signal (as point-biserial $r = -0.022$ also says).

\textbf{R18. Group means and a Welch t-test.}
\begin{lstlisting}
a = df.loc[df.target == 1, "city_development_index"]; b = df.loc[df.target == 0, "city_development_index"]
a.mean(), b.mean(), stats.ttest_ind(a, b, equal_var=False).pvalue      # 0.756, 0.853, ~0
\end{lstlisting}
\textbf{R19. Chi-square for a categorical feature:}
\begin{lstlisting}
stats.chi2_contingency(pd.crosstab(df["company_type"].fillna("Unknown"), df["target"]))[1]   # 4e-204
\end{lstlisting}
With 19{,}000 rows almost anything is ``significant''; use Cram\'er's $V$ for strength.

\textbf{R20. Pearson vs Spearman with the target.}
\begin{lstlisting}
num = df[["city_development_index", "training_hours", "exp_num", "target"]]
num.corr()["target"], num.corr(method="spearman")["target"]
\end{lstlisting}
\outlabel
\begin{lstlisting}[style=out]
pearson : CDI -0.342  training_hours -0.022  exp_num -0.177
spearman: CDI -0.279  training_hours -0.014  exp_num -0.184
\end{lstlisting}

\subsection{Step 11: interactions and segments}
\textbf{R17. Two-factor pivot.}
\begin{lstlisting}
df.pivot_table(index="education_level", columns="relevent_experience", values="target", aggfunc="mean")
\end{lstlisting}
\outlabel
\begin{lstlisting}[style=out]
                 Has relevent experience  No relevent experience
Graduate                          0.23                    0.43
High School                       0.12                    0.24
Masters                           0.19                    0.32
Phd                               0.12                    0.17
Primary School                    0.15                    0.13
\end{lstlisting}
Graduates without relevant experience are the most mobile group (43\%): a candidate segment for targeted job alerts.

\subsection{Feature-engineering recipes that come out of EDA}
\textbf{R21. Relative-to-group feature with \code{transform}:}
\begin{lstlisting}
df["exp_vs_city_avg"] = df["exp_num"] - df.groupby("city")["exp_num"].transform("mean")
\end{lstlisting}
\textbf{R22. Frequency encoding:}
\begin{lstlisting}
df["city_freq"] = df["city"].map(df["city"].value_counts(normalize=True))     # city_103 -> 0.2273
\end{lstlisting}
\textbf{R24. Multi-condition labels with \code{np.select}:}
\begin{lstlisting}
df["career_stage"] = np.select([df.exp_num <= 2, df.exp_num <= 10], ["junior", "mid"], default="senior")
df.groupby("career_stage")["target"].mean()        # junior 0.384, mid 0.284, senior 0.171
\end{lstlisting}

\section{Plot recipes (Matplotlib/Seaborn)}
\begin{lstlisting}
fig, axes = plt.subplots(1, 4, figsize=(15, 3))                       # grid of histograms
for ax, c in zip(axes, num_cols): sns.histplot(df[c], kde=True, ax=ax); ax.set_title(c)
sns.boxplot(data=df, x="target", y="city_development_index")          # numeric by class
sns.kdeplot(data=df, x="city_development_index", hue="target", common_norm=False)  # overlap
rate_by(df, "company_type", "target")["rate"].plot(kind="barh")       # rate by category
sns.heatmap(df[num_cols].corr(), annot=True, fmt=".2f", cmap="coolwarm", vmin=-1, vmax=1)
sns.pointplot(data=df, x="career_stage", y="target", hue="relevent_experience", errorbar=None)  # interaction
sns.scatterplot(data=df, x="x1", y="x2", hue="target", alpha=0.5)    # two features at once
\end{lstlisting}
Rules: \code{common\_norm=False} when comparing distributions of unequal-sized classes; always show $n$ beside rates; sort bar charts; log axes for skewed money.

\section{How to narrate EDA in the interview}
\begin{enumerate}
\item Start with the frame: ``One row is a candidate; target is job change; 25\% positive, so I will use PR AUC/F1 and stratified splits.''
\item Data quality before statistics: types, codes, duplicates, malformed labels.
\item Missingness with the target, not just counts.
\item Rank features by separation (point-biserial, Cram\'er's $V$, binned rates) and say which you would expect from domain sense.
\item Name one interaction and one non-linearity.
\item Close with a written plan: drops (leaks, IDs), encodings, transforms, imputation, model choice, metric.
\end{enumerate}

\stuck{pandas User Guide: Group by (split-apply-combine)}{https://pandas.pydata.org/docs/user_guide/groupby.html}
\section{Interview questions}

\begin{qa}{\pyq{Round 3 Part 1} You get a new CSV in a Colab notebook. What are your first ten lines of code, and why each?}
\oneline{Load, then shape/head/info, the overview table (types, missing \%, unique counts), duplicates, target balance, \code{describe} with skew and zero counts, value counts of
categoricals, missingness vs target, and a quick leakage check of any column that perfectly separates the target.}
\why Each line answers a question that changes a decision (Section 1). Leakage and data quality first, statistics second, plots third, model last.
\worked
\begin{lstlisting}
df = pd.read_csv(path); print(df.shape); df.head()
overview(df)
df.duplicated().sum(), df.drop(columns=id_col).duplicated().sum()
df[target].value_counts(normalize=True)
num_summary(df)
for c in df.select_dtypes(exclude="number"): print(df[c].value_counts(dropna=False).head(10))
missing_vs_target(df, target)
point_biserial(df, num_cols, target)
\end{lstlisting}
\end{qa}

\begin{qa}{\deep \code{isna().sum()} shows zero missing values. Are you done with missing values?}
\oneline{No: missing values are often disguised as 0, $-1$, 999, empty strings, ``NA'' text or a default date; check zero counts, impossible values, spikes at round numbers and
value counts of text columns.}
\worked The Pima diabetes data (Chapter 44) has no NaN but 374 zero Insulin and 35 zero blood-pressure readings: impossible values that are really missing.
\end{qa}

\begin{qa}{How do you find out whether missing values are informative?}
\oneline{Compare the target rate (or target mean) for rows where the field is missing vs present, and check what else is missing with it; a big difference means missingness carries
signal and should be encoded (indicator or ``Unknown''), not imputed away.}
\worked HR data: job-change rate 0.406 when company size is blank vs 0.179 when present; company size and type are missing together in 5{,}360 rows.
\end{qa}

\begin{qa}{Why do you print counts next to rates in every group-by?}
\oneline{Because a rate from a tiny group is unreliable; 100\% from 2 people is noise, while 40\% from 6{,}000 is solid; counts tell the listener how much to trust each number and
whether to group rare levels.}
\end{qa}

\begin{qa}{\deep Training hours shows 984 IQR outliers and 450 z-score outliers. Do you remove them?}
\oneline{No: the variable is right-skewed, so both rules flag many genuine heavy learners; they are valid values, so I keep them and, for linear models, log-transform or cap; tree
models are unaffected.}
\why Outlier rules assume a roughly symmetric distribution; on skewed data they flag the long tail by design. Remove only impossible values.
\end{qa}

\begin{qa}{How do you turn EDA findings into a modelling plan? Give the plan for the HR data.}
\oneline{Translate each finding into an action: drop the ID; parse experience and last job; fix the \code{10/49} label; keep company missingness as a category; frequency- or
target-encode city; keep CDI (strongest, non-linear, so a tree model); stratified split; class weights; evaluate with PR AUC and F1.}
\end{qa}

\begin{qa}{What is the difference between \code{groupby().agg()}, \code{transform()} and \code{pivot\_table()} in EDA?}
\oneline{\code{agg} gives one summary row per group (rates by city); \code{transform} returns a group statistic aligned to every original row (features relative to the group);
\code{pivot\_table} lays out a statistic across two grouping variables (interactions).}
\end{qa}

\begin{qa}{\deep Your target rate by a feature looks almost perfect (0\% vs 100\%). What do you suspect and do?}
\oneline{Leakage: the feature is probably created by or after the outcome (salary for placement, ``interview scheduled'' for hiring); confirm with domain questions and timing, then drop it
from the predictive model (it may still be useful for a different question).}
\end{qa}
